% !TeX program = xelatex

\documentclass[UTF8,zihao=-4]{ctexart}

\usepackage[a4paper,margin=2.3cm]{geometry}
\usepackage{amsmath,amssymb}
\usepackage{booktabs}
\usepackage{longtable}
\usepackage{graphicx}
\usepackage{float}
\usepackage{array}
\usepackage{xcolor}
\usepackage{hyperref}
\usepackage{caption}
\usepackage{enumitem}

\hypersetup{
    colorlinks=true,
    linkcolor=blue,
    urlcolor=blue,
    citecolor=blue
}

\setlist[itemize]{leftmargin=2em}
\setlist[enumerate]{leftmargin=2em}
\captionsetup{font=small}

% 如果图片文件尚未上传，报告仍可正常编译，并显示图片占位框。
% 上传同名图片后重新编译，即可自动替换为真实截图。
\newcommand{\maybeincludegraphics}[2][]{%
  \IfFileExists{#2}{\includegraphics[#1]{#2}}{%
    \fbox{\begin{minipage}[c][4.0cm][c]{0.88\linewidth}\centering
    图片占位：请上传 \texttt{\detokenize{#2}}\par
    上传后重新编译即可显示
    \end{minipage}}%
  }%
}


\title{\textbf{Warping 后门攻击的复现与频域检测分析}\\
\large WaNet: Imperceptible Warping-based Backdoor Attack 课程报告}
\author{周争妍 \quad 蒋竺君\\
人工智能安全技术课程设计}
\date{\today}

\begin{document}
\maketitle

\begin{abstract}
后门攻击使深度神经网络在正常样本上保持较高准确率，但在带有特定触发器的输入上输出攻击者指定标签。传统后门触发器通常表现为可见贴片、水印或局部扰动，因此较容易被人工观察或现有防御方法发现。WaNet 提出使用图像空间形变作为触发器，通过微小且人眼难以察觉的 warping 变换建立输入与目标标签之间的关联。

本报告围绕 WaNet 进行了复现、消融与扩展分析。首先，复现了 CIFAR-10 上的主要攻击效果，并验证 STRIP、Fine-Pruning 和 Neural Cleanse 等防御方法难以稳定检测或移除该后门。其次，通过 Noise Mode、形变强度 $s$ 与控制网格大小 $k$ 的消融实验，分析 WaNet 成功的关键机制。最后，提出一个基于频域特征的轻量检测器：将 clean 图像与 warped 图像转到傅里叶域，提取高频能量占比、频谱熵等 6 维特征，并使用 RBF-SVM 进行分类。实验结果表明，该频域检测器在 CIFAR-10 上达到 AUC=0.735，相比 STRIP 近似随机猜测的 AUC=0.50 有明显提升。
\end{abstract}

\section{引言}

深度学习模型在图像分类、目标识别等任务中已经取得很高性能，但模型训练流程往往依赖外部数据、预训练模型或第三方训练服务，这使得后门攻击成为实际安全风险。后门攻击的核心目标是：模型在干净输入上表现正常，但当输入中出现攻击者指定的触发器时，模型输出被攻击者控制的目标类别。

传统后门攻击常用触发器包括固定白色方块、局部水印、噪声图案或正弦纹理。这些方法实现简单，但触发器通常具有人眼可见的局部异常，也更容易被 STRIP、Neural Cleanse、Fine-Pruning 等防御方法利用输出熵、逆向触发器搜索或神经元剪枝发现。WaNet 的关键创新在于：不再向图像叠加显式图案，而是对整张图像做微小空间形变，使触发器以几何变换的形式存在。

本项目主要完成三部分工作：
\begin{enumerate}
    \item \textbf{主实验复现：}复现 WaNet 在 CIFAR-10 上的攻击效果，并比较 Clean Accuracy、Attack Success Rate 和 Cross Accuracy。
    \item \textbf{防御与消融分析：}验证 STRIP、Fine-Pruning、Neural Cleanse 对 WaNet 的失效现象，并分析 Noise Mode、形变强度 $s$、控制网格大小 $k$ 对攻击效果的影响。
    \item \textbf{扩展研究：}从图像信号本身出发，设计频域 warping 检测器，验证 warping 触发器虽然人眼不可察觉，但仍可能在频域留下统计痕迹。
\end{enumerate}

\section{相关工作与方法背景}

\subsection{后门攻击问题定义}

设分类模型为 $f_\theta(x)$，干净训练集为 $\mathcal{D}=\{(x_i,y_i)\}_{i=1}^n$。后门攻击者在训练阶段构造一部分带触发器的样本：
\[
    x_i' = T(x_i), \qquad y_i' = y_t,
\]
其中 $T(\cdot)$ 为触发器注入函数，$y_t$ 为目标标签。训练后的模型需要同时满足：
\[
    f_\theta(x) \approx y \quad \text{for clean samples},
\]
\[
    f_\theta(T(x)) \approx y_t \quad \text{for triggered samples}.
\]
因此，一个成功的后门攻击既要保持干净样本准确率，又要使触发样本具有高攻击成功率。

\subsection{WaNet 的核心思想}

WaNet 使用图像空间形变代替可见贴片。其做法是先生成一个低分辨率随机控制网格，再通过插值得到平滑位移场，最后对原图进行轻微 warping。可以将该过程抽象为：
\[
    x' = W(x; G),
\]
其中 $W$ 表示空间采样函数，$G$ 是形变网格。实际实现中，WaNet 使用类似 PyTorch \texttt{grid\_sample} 的双线性插值操作，对图像坐标进行亚像素级扰动。由于这种变换不引入明显局部图案，触发器在视觉上更加隐蔽。

一个简化的形变网格可以表示为：
\[
    G = I + \frac{s}{H}\Delta,
\]
其中 $I$ 是 identity grid，$\Delta$ 是由随机控制网格上采样得到的平滑位移场，$s$ 控制形变强度，$H$ 是图像高度。控制网格大小 $k$ 决定形变模式的空间频率：$k$ 太小会导致形变过于粗糙，$k$ 太大则可能使形变过细或接近噪声。

\subsection{Noise Mode}

WaNet 的另一个关键设计是 Noise Mode。训练时除了加入真正的后门样本，还加入带随机形变但标签不改变的 noise samples。这样做的目的不是增强攻击强度，而是避免模型学习到简单的像素级捷径。换句话说，Noise Mode 迫使模型区分"特定 warping 模式"和"一般随机形变"，从而使后门触发器更难被 Neural Cleanse 等方法逆向恢复。

\section{实验设置}

\subsection{数据集与模型}

本项目的实验分为两层：第一层是主复现实验，覆盖 MNIST、CIFAR-10、GTSRB 三个数据集，以及 all2one、all2all 两种攻击模式，用来验证 WaNet 在不同数据和攻击设置下的整体可复现性；第二层是后续分析实验，重点在 CIFAR-10 上展开 Noise Mode、形变强度 $s$、控制网格大小 $k$ 的消融分析和频域检测器扩展。复现实验采用 WaNet 官方实现中的训练与评估流程，并结合本项目仓库中的脚本完成环境适配、结果保存、防御实验运行与截图整理。

\subsection{评价指标}

本文使用以下指标评价攻击与检测效果：
\begin{itemize}
    \item \textbf{Clean Acc：}模型在干净测试集上的分类准确率。该指标越高，说明后门对正常性能破坏越小。
    \item \textbf{ASR：}Attack Success Rate，即带触发器样本被分类到攻击目标类别的比例。该指标越高，说明攻击越有效。
    \item \textbf{Cross Acc：}使用随机形变或 noise samples 进行测试时的准确率。该指标用于衡量模型是否真正学习了特定 warping，而不是对任意形变都误触发。
    \item \textbf{AUC：}频域检测器的 ROC 曲线下面积。AUC=0.5 表示近似随机猜测，AUC 越高表示检测能力越强。
\end{itemize}

\subsection{防御方法}

我们验证了三类经典后门防御方法：
\begin{itemize}
    \item \textbf{STRIP：}通过对输入进行扰动并观察预测熵的变化来判断样本是否带触发器。如果触发器强制模型输出固定标签，则扰动后的输出熵通常较低。
    \item \textbf{Fine-Pruning：}剪除在干净样本上激活较弱的神经元，希望移除后门相关神经元，同时尽量保留正常分类能力。
    \item \textbf{Neural Cleanse：}对每个类别反向优化可能的最小触发器，并通过 anomaly index 判断是否存在异常小的目标类别触发器。
\end{itemize}


\section{复现工作流}

本节单独展开主复现部分。与后文的 CIFAR-10 消融实验不同，主复现实验的重点不是只跑一个 baseline，而是完成环境配置、数据下载、训练评估、三类防御验证、日志管理和结果截图整理，覆盖 3 个数据集与 2 种攻击模式。

\subsection{复现范围}

主复现部分覆盖的数据集、攻击模式和防御方法如表 \ref{tab:teammate_scope} 所示。

\begin{table}[H]
\centering
\caption{主复现范围}
\label{tab:teammate_scope}
\begin{tabular}{p{0.22\linewidth}p{0.68\linewidth}}
\toprule
模块 & 具体内容 \\
\midrule
数据集 & MNIST、CIFAR-10、GTSRB \\
攻击模式 & all2one、all2all，共 $3\times2=6$ 组主实验 \\
主指标 & Clean test accuracy、Attack test accuracy、Noise/Cross test accuracy \\
防御方法 & STRIP、Fine-Pruning、Neural Cleanse \\
运行方式 & tmux 后台任务；GPU0/GPU1 并行；日志文件与 summary.tsv 汇总 \\
交付材料 & 训练/评估脚本、结果 JSON、STRIP/FP/NC 输出、tmux 截图、日志截图、结果截图 \\
\bottomrule
\end{tabular}
\end{table}

\subsection{运行环境与工程适配}

复现实验不是直接运行单个 notebook，而是在服务器环境中组织为可重复执行的脚本工作流。主要环境配置如下：

\begin{itemize}
    \item Python 3.10；
    \item PyTorch 2.6.0 + CUDA 12.4；
    \item torchvision 0.21.0；
    \item 默认使用 \texttt{CUDA\_VISIBLE\_DEVICES=0}，并通过 tmux 将不同实验分配到后台 worker；
    \item 主要代码入口包括 \texttt{train.py}、\texttt{eval.py}、\texttt{defenses/STRIP}、\texttt{defenses/fine\_pruning}、\texttt{defenses/neural\_cleanse}。
\end{itemize}

主实验运行流程可以概括为：

\begin{enumerate}
    \item 下载并检查 CIFAR-10 / MNIST / GTSRB 数据集；
    \item 对每个数据集分别运行 all2one 和 all2all 训练；
    \item 对训练后 checkpoint 执行 clean / attack / noise 三类测试；
    \item 对每组模型运行 STRIP 和 Fine-Pruning；
    \item 对 MNIST all2one、CIFAR-10 all2one 补充运行 Neural Cleanse；
    \item 汇总 \texttt{eval\_results.json}、防御输出文件和 \texttt{summary.tsv}，并截图保存到 \texttt{report\_assets/}。
\end{enumerate}

核心运行命令示例如下：

\begin{verbatim}
cd /home/zilia/project/experiment/wanet
bash scripts/start_weekend_tmux.sh
bash scripts/status_weekend.sh
tail -f logs/weekend_latest/summary.tsv
\end{verbatim}

单独的 CIFAR-10 all2one 训练与评估命令如下：

\begin{verbatim}
bash scripts/train_cifar10_all2one.sh
bash scripts/eval_cifar10_all2one.sh
\end{verbatim}

防御实验命令示例如下：

\begin{verbatim}
bash scripts/strip_cifar10_all2one.sh
cd defenses/fine_pruning
python fine-pruning-cifar10-gtsrb.py --dataset cifar10 --attack_mode all2one
cd ../neural_cleanse
python neural_cleanse.py --dataset cifar10 --attack_mode all2one
\end{verbatim}

\subsection{工作流与截图材料}

为了让报告更直观地体现实际工作量，建议在本节加入真实运行截图。下面的图片文件名已经写入 LaTeX；你只需要在 Overleaf 左侧上传同名图片，重新编译后占位框会自动替换为真实图片。

\begin{figure}[H]
\centering
\maybeincludegraphics[width=0.92\linewidth]{wanet_paper_teaser.png}
\caption{WaNet 原论文中的触发器对比图。相比 Patched、Blended、SIG 和 ReFool 等可见或半可见触发器，WaNet 使用空间形变作为触发器，在视觉上更难察觉。图源：Nguyen \& Tran, ICLR 2021。}
\label{fig:paper_teaser}
\end{figure}

\begin{figure}[H]
\centering
\maybeincludegraphics[width=0.92\linewidth]{tmux_gpu0_final.png}
\caption{GPU0 worker 的 tmux 最终运行界面截图，用于展示后台训练与评估任务完成情况。}
\label{fig:tmux_gpu0}
\end{figure}

\begin{figure}[H]
\centering
\maybeincludegraphics[width=0.92\linewidth]{tmux_gpu1_final.png}
\caption{GPU1 worker 的 tmux 最终运行界面截图，用于展示并行任务调度和多实验运行情况。}
\label{fig:tmux_gpu1}
\end{figure}

\begin{figure}[H]
\centering
\maybeincludegraphics[width=0.92\linewidth]{summary_tsv_tail.png}
\caption{\texttt{summary.tsv} 运行日志尾部截图，展示训练、评估和防御任务的完成记录。}
\label{fig:summary_tsv}
\end{figure}

\begin{figure}[H]
\centering
\maybeincludegraphics[width=0.92\linewidth]{eval_results_json.png}
\caption{\texttt{eval\_results.json} 结果文件截图，展示 clean、attack、noise/cross 三类指标的真实输出。}
\label{fig:eval_json}
\end{figure}

\begin{figure}[H]
\centering
\maybeincludegraphics[width=0.92\linewidth]{strip_results.png}
\caption{STRIP 防御输出文件截图，展示 trojan entropy 与 benign entropy 的统计结果。}
\label{fig:strip_results}
\end{figure}

\begin{figure}[H]
\centering
\maybeincludegraphics[width=0.92\linewidth]{fine_pruning_results.png}
\caption{Fine-Pruning 防御结果截图，展示剪枝过程中 clean accuracy 与 attack accuracy 的变化。}
\label{fig:fp_results}
\end{figure}

\begin{figure}[H]
\centering
\maybeincludegraphics[width=0.92\linewidth]{neural_cleanse_results.png}
\caption{Neural Cleanse 输出截图，展示 anomaly index 和逆向触发器检测结果。}
\label{fig:nc_results}
\end{figure}

\subsection{主复现实验结果}

表 \ref{tab:teammate_main_results} 给出了 3 个数据集、2 种攻击模式下的主复现结果。可以看到，6 个设置均保持了较高的 clean accuracy 和 attack accuracy，说明复现实验没有只依赖单个数据集或单个攻击模式。

\begin{table}[H]
\centering
\caption{3 数据集 $\times$ 2 攻击模式主复现结果}
\label{tab:teammate_main_results}
\resizebox{\linewidth}{!}{%
\begin{tabular}{llcccccc}
\toprule
数据集 & 攻击模式 & 复现 Clean & 复现 Attack & 复现 Noise & 参考 Clean & 参考 Attack & 参考 Noise \\
\midrule
MNIST & all2one & 99.50 & 99.98 & 99.16 & 99.52 & 99.86 & 98.20 \\
CIFAR-10 & all2one & 94.52 & 99.37 & 92.52 & 94.15 & 99.55 & 93.55 \\
GTSRB & all2one & 99.34 & 99.75 & 99.03 & 98.87 & 99.33 & 98.01 \\
MNIST & all2all & 99.45 & 99.16 & 99.18 & 99.44 & 95.90 & 94.34 \\
CIFAR-10 & all2all & 94.36 & 93.52 & 92.65 & 94.43 & 93.36 & 91.47 \\
GTSRB & all2all & 99.59 & 99.22 & 99.14 & 99.39 & 98.31 & 98.96 \\
\bottomrule
\end{tabular}}
\end{table}

从表中可以看到，CIFAR-10 all2one 的复现 Clean Acc 为 94.52\%，Attack Acc 为 99.37\%，与参考结果 94.15\%、99.55\% 基本一致。GTSRB 和 MNIST 的结果也接近或高于参考值，说明主复现流程整体稳定。需要注意的是，本节结果来自完整主复现实验；后文 CIFAR-10 消融部分使用的是另一组 baseline 记录，因此数值存在小幅差异是正常的。

\subsection{防御验证结果}

表 \ref{tab:strip_main} 展示了 STRIP 的关键输出。STRIP 的核心假设是：带触发器样本在被叠加其他图像扰动后，模型仍会稳定输出目标类，因此预测熵会显著偏低。但在本次 WaNet 复现中，trojan entropy 并没有稳定低到可检测阈值以下，说明 STRIP 对 warping trigger 区分能力较弱。

\begin{table}[H]
\centering
\caption{STRIP 防御结果摘要}
\label{tab:strip_main}
\resizebox{\linewidth}{!}{%
\begin{tabular}{llccc}
\toprule
数据集 & 攻击模式 & Trojan entropy 最小值 & Trojan entropy 均值 & Benign entropy 均值 \\
\midrule
MNIST & all2one & 1.177 & 1.703 & 1.772 \\
MNIST & all2all & 1.734 & 2.295 & 1.916 \\
CIFAR-10 & all2one & 2.216 & 2.684 & 2.637 \\
CIFAR-10 & all2all & 2.019 & 2.761 & 2.700 \\
GTSRB & all2one & 9.270 & 11.999 & 10.588 \\
GTSRB & all2all & 10.189 & 12.804 & 12.565 \\
\bottomrule
\end{tabular}}
\end{table}

Fine-Pruning 的结果也显示，当 clean accuracy 仍保持在 90\% 以上时，attack accuracy 往往仍然很高。例如 CIFAR-10 all2one 在剪枝后 clean accuracy 约为 90.02\% 时，attack accuracy 仍约为 99.91\%。这说明 WaNet 后门并不容易通过简单神经元剪枝移除。

\begin{table}[H]
\centering
\caption{Fine-Pruning 防御结果摘要}
\label{tab:fp_main}
\resizebox{\linewidth}{!}{%
\begin{tabular}{llcc}
\toprule
数据集 & 攻击模式 & 原始 Clean / Attack & Clean $>$ 90\% 时最后结果 \\
\midrule
MNIST & all2one & 99.50 / 99.98 & 90.69 / 99.91 \\
MNIST & all2all & 99.45 / 99.16 & 90.22 / 95.97 \\
CIFAR-10 & all2one & 94.52 / 99.37 & 90.02 / 99.91 \\
CIFAR-10 & all2all & 94.36 / 93.52 & 90.36 / 86.74 \\
GTSRB & all2one & 99.34 / 99.75 & 90.27 / 99.76 \\
GTSRB & all2all & 99.59 / 99.22 & 90.35 / 70.38 \\
\bottomrule
\end{tabular}}
\end{table}

Neural Cleanse 方面，主复现中完成了 MNIST all2one 和 CIFAR-10 all2one 两组实验。由于 Neural Cleanse 本身对初始化、优化过程和触发器搜索较敏感，本次实验没有得到稳定强异常信号，整体符合 WaNet 难以被 Neural Cleanse 稳定检测的趋势。

\subsection{主复现部分小结}

周争妍负责的部分主要体现为完整复现链路，而不是单个数值结果：从环境配置、数据下载、训练脚本、tmux 并行管理，到 6 组主实验、STRIP/Fine-Pruning/Neural Cleanse 防御验证和结果截图整理。这部分工作证明了 WaNet 的攻击效果与防御绕过现象可以在多个数据集和攻击模式下复现，为后续消融分析和频域检测扩展提供了可信 baseline。

\section{主实验结果}

\subsection{CIFAR-10 主复现结果}

表 \ref{tab:main} 展示了 CIFAR-10 上的主实验结果。可以看到，本次复现得到 Clean Acc=92.84\%、ASR=99.15\%，说明模型在保持较高正常分类能力的同时，仍能在触发样本上几乎稳定地输出目标标签。与论文参考结果相比，Clean Acc 和 Cross Acc 略低，但 ASR 基本一致，说明攻击效果已经成功复现。

\begin{table}[H]
\centering
\caption{CIFAR-10 主实验复现结果}
\label{tab:main}
\begin{tabular}{lccc}
\toprule
设置 & Clean Acc (\%) & ASR (\%) & Cross Acc (\%) \\
\midrule
本项目复现 & 92.84 & 99.15 & 89.71 \\
论文参考结果 & 94.15 & 99.55 & 93.55 \\
\bottomrule
\end{tabular}
\end{table}

从结果看，WaNet 的攻击成功率非常高，而干净准确率没有崩溃。这正是后门攻击最危险的地方：如果只看普通测试集准确率，模型看起来仍然是一个正常模型。

\subsection{防御实验结果}

防御实验结论如下：
\begin{itemize}
    \item \textbf{STRIP：}clean 样本与 trojan 样本的 entropy 分布难以区分，检测效果接近随机。
    \item \textbf{Fine-Pruning：}在 Clean Acc 仍保持较高水平时，ASR 仍然大于 90\%，说明简单剪枝不能有效移除 WaNet 后门。
    \item \textbf{Neural Cleanse：}在开启 Noise Mode 的设置下，Anomaly Index 小于常用阈值 2，未能稳定检测到异常目标类。
\end{itemize}

这些结果说明 WaNet 的隐蔽性不只体现在视觉层面，也体现在模型行为层面。STRIP 依赖输出熵，Neural Cleanse 依赖可逆向恢复的小触发器，Fine-Pruning 依赖后门神经元可分离性；而 WaNet 的触发器是全局形变，且在 Noise Mode 训练下不再表现为简单局部模式，因此上述防御均会失效。

\section{消融实验分析}

\subsection{Noise Mode 的作用}

表 \ref{tab:noise} 展示了有无 Noise Mode 的对比。可以看到，关闭 Noise Mode 后，Clean Acc 和 ASR 仍然较高，但 Cross Acc 降为 0。这说明模型没有学到真正稳健的 warping 模式，而是学到了某种容易被检测的像素捷径。与此同时，Neural Cleanse 在无 Noise Mode 设置下可以检测到异常。

\begin{table}[H]
\centering
\caption{Noise Mode 消融实验}
\label{tab:noise}
\begin{tabular}{lcccc}
\toprule
设置 & Clean Acc (\%) & ASR (\%) & Cross Acc (\%) & Neural Cleanse \\
\midrule
With Noise Mode & 92.84 & 99.15 & 89.71 & 检测不到 \\
Without Noise Mode & 92.74 & 98.21 & 0.00 & 能检测到 \\
\bottomrule
\end{tabular}
\end{table}

这组结果是理解 WaNet 的关键。Noise Mode 并不是简单的数据增强，而是 WaNet 绕过防御的核心训练机制。没有 Noise Mode 时，攻击仍然可能成功，但隐蔽性明显下降。

\subsection{形变强度 \texorpdfstring{$s$}{s} 的影响}

表 \ref{tab:s} 展示了不同形变强度 $s$ 下的结果。$s=0.25$ 时 ASR 仅为 74.12\%，说明形变过弱会导致触发器信号不足；$s=0.50$ 时 ASR 达到 99.59\%，攻击最稳定；$s=1.00$ 时 ASR 仍然很高，但收益不再明显增加。

\begin{table}[H]
\centering
\caption{Warping 强度 $s$ 的消融实验}
\label{tab:s}
\begin{tabular}{lccc}
\toprule
$s$ & Clean Acc (\%) & ASR (\%) & Cross Acc (\%) \\
\midrule
0.25 & 92.39 & 74.12 & 89.38 \\
0.50 & 92.81 & 99.59 & 89.69 \\
1.00 & 92.82 & 98.29 & 89.30 \\
\bottomrule
\end{tabular}
\end{table}

因此，WaNet 并不是任意微小形变都能成功。攻击存在一个有效参数区间：形变太弱，模型难以学习；形变太强，虽然攻击仍可成功，但视觉隐蔽性可能下降。

\subsection{控制网格大小 \texorpdfstring{$k$}{k} 的影响}

表 \ref{tab:k} 展示了不同控制网格大小 $k$ 下的结果。$k=2$ 时 ASR 为 92.82\%，Cross Acc 为 88.95\%，均明显低于 $k=4$；$k=4$ 时 ASR 达到 99.15\%；$k=8$ 时 ASR 和 Cross Acc 仍保持较高水平，但不再进一步提升。

\begin{table}[H]
\centering
\caption{控制网格大小 $k$ 的消融实验}
\label{tab:k}
\begin{tabular}{lcccc}
\toprule
$k$ & Epochs & Clean Acc (\%) & ASR (\%) & Cross Acc (\%) \\
\midrule
2 & 80 & 92.14 & 92.82 & 88.95 \\
4 & 400 & 92.84 & 99.15 & 89.71 \\
8 & 80 & 91.94 & 97.55 & 88.04 \\
\bottomrule
\end{tabular}
\end{table}

该结果说明，warping 模式的空间尺度很重要。控制网格太粗时，形变模式表达能力不足；控制网格增大后，触发器更容易被模型学习，但过大后边际收益下降。需注意 $k=2$ 和 $k=8$ 仅训练了 80 epochs（vs $k=4$ 的 400 epochs），因此绝对数值偏低，但 ASR 的相对趋势仍然有效。

\subsection{多种子稳定性验证}

为验证结果的可复现性，在默认参数 ($s=0.5$, $k=4$) 下使用不同随机种子（0 和 1）训练两次，结果如表 \ref{tab:seed} 所示。两组实验的 Clean Acc、ASR 和 Cross Acc 均保持在相近水平，表明 WaNet 的攻击效果在相同参数设定下具有统计一致性。

\begin{table}[H]
\centering
\caption{多种子稳定性验证}
\label{tab:seed}
\begin{tabular}{cccc}
\toprule
Seed & Clean Acc (\%) & ASR (\%) & Cross Acc (\%) \\
\midrule
0 & 92.81 & 99.59 & 89.69 \\
1 & 92.10 & 98.70 & 88.72 \\
\bottomrule
\end{tabular}
\end{table}

\section{频域检测器}

\subsection{动机}

虽然 WaNet 的触发器在人眼看来几乎不可见，但其实现依赖空间采样和双线性插值。插值本质上是一种信号处理操作：新像素由周围像素加权平均得到。这种操作可能改变图像的频率分布，尤其是在高频能量、频谱熵和径向频率统计上留下痕迹。

现有防御方法大多观察模型行为，例如输出熵、神经元激活或逆向触发器大小。我们的思路不同：直接观察输入图像信号本身。如果 warping 触发器改变了图像频域统计，那么即使模型输出行为无法区分，图像层面仍可能提供检测线索。

\subsection{方法}

频域检测器流程如下：
\begin{enumerate}
    \item 构造 clean 图像与 warped 图像各 500 张。
    \item 对每张图像做二维快速傅里叶变换，并将低频移动到频谱中心。
    \item 对每个通道分别提取频域特征，再对通道求平均。
    \item 使用 6 维频域特征训练 RBF 核 SVM。
    \item 使用 30\% 测试集计算准确率与 ROC AUC。
\end{enumerate}

提取的 6 维特征为：
\begin{itemize}
    \item 高频能量占比；
    \item 频谱熵；
    \item 频率加权均值；
    \item 频率加权方差；
    \item 低频能量占比；
    \item 中频能量占比。
\end{itemize}

以单通道图像 $x$ 为例，设其傅里叶幅度谱为 $M(u,v)$，总能量为
\[
    E = \sum_{u,v} M(u,v).
\]
高频能量占比可写为
\[
    R_{\text{high}} = \frac{\sum_{(u,v)\in \Omega_{\text{high}}} M(u,v)}{E}.
\]
频谱熵定义为
\[
    H = -\sum_{u,v} p(u,v)\log_2 p(u,v), \qquad
    p(u,v)=\frac{M(u,v)}{E}.
\]
这些统计量不依赖模型输出，而是直接刻画图像本身的频率结构。

\subsection{频域可视化}

图 \ref{fig:fft} 展示了 Clean 图像与 Warped 图像的 FFT 频谱对比。可以观察到，warped 图像的频谱在高频区域存在可辨识的结构化差异——这是插值平滑操作在频域留下的痕迹。右下角的频谱差异图更直观地展示了这种异常：差异并非均匀分布，而是集中在特定频率带上，与插值对高频分量的改变一致。

\begin{figure}[H]
\centering
\maybeincludegraphics[width=0.88\linewidth]{fft_spectrum_comparison.jpg}
\caption{Clean 图像与 Warped 图像的 FFT 频谱对比。上排：Clean 图像及其 FFT 频谱；下排：Warped 图像及其 FFT 频谱；右下角：频谱差异图（红色区域表示差异显著）。}
\label{fig:fft}
\end{figure}

图 \ref{fig:feat_dist} 进一步展示了六维频域特征在 clean 和 warped 样本上的分布直方图。高频能量比、频谱熵和方差频率在两类样本上存在明显的分布分离，从统计层面验证了 warping 对图像频域结构的实质性影响。

\begin{figure}[H]
\centering
\maybeincludegraphics[width=\linewidth]{feature_distributions.jpg}
\caption{六维频域特征在 Clean（蓝）和 Warped（红）样本上的分布。可观察到高频能量比、频谱熵等维度上两类样本分布有显著差异。}
\label{fig:feat_dist}
\end{figure}

\subsection{检测结果}

表 \ref{tab:freq} 展示了频域检测器与 STRIP 的对比。频域检测器达到 AUC=0.735，准确率为 67.33\%；而 STRIP 的 AUC 约为 0.50，基本等同随机猜测。

\begin{table}[H]
\centering
\caption{频域检测器与 STRIP 对比}
\label{tab:freq}
\begin{tabular}{lcc}
\toprule
方法 & Accuracy (\%) & AUC \\
\midrule
频域检测器（Ours） & 67.33 & 0.735 \\
STRIP & -- & 0.50 \\
\bottomrule
\end{tabular}
\end{table}

图 \ref{fig:roc} 展示了频域检测器与 STRIP 的 ROC 曲线对比。蓝色实线为频域检测器，AUC 达到 0.735；红色标注点为 STRIP，AUC 约为 0.50（即沿对角线的随机猜测水平）。两者差异直观地说明了频域视角的检测优势。

\begin{figure}[H]
\centering
\maybeincludegraphics[width=0.58\linewidth]{roc_curve.jpg}
\caption{频域检测器与 STRIP 的 ROC 曲线对比。蓝色：频域检测器（AUC=0.735）；红色X：STRIP（AUC=0.50）。}
\label{fig:roc}
\end{figure}

这个结果不能说明频域检测器已经足够作为强防御系统，因为 67.33\% 的准确率距离实际部署仍有差距；但它说明一个重要事实：WaNet 的"不可感知"并不等于"不可检测"。模型行为层面的防御失败后，图像信号层面仍然可能存在可利用信息。

\section{总体结果汇总}

表 \ref{tab:summary} 汇总了本项目主要实验结果。

\begin{longtable}{lcccl}
\caption{实验结果汇总}
\label{tab:summary}\\
\toprule
实验设置 & Epochs & Clean Acc (\%) & ASR (\%) & 关键结论 \\
\midrule
\endfirsthead
\toprule
实验设置 & Epochs & Clean Acc (\%) & ASR (\%) & 关键结论 \\
\midrule
\endhead
Noise WITH & 400 & 92.84 & 99.15 & Baseline，攻击有效 \\
Noise WITHOUT & 400 & 92.74 & 98.21 & Cross Acc 为 0，隐蔽性下降 \\
$s=0.25$ & 400 & 92.39 & 74.12 & 形变过弱，攻击明显下降 \\
$s=0.50$ & 400 & 92.81 & 99.59 & 最优强度附近 \\
$s=1.00$ & 400 & 92.82 & 98.29 & 攻击饱和 \\
$k=2$ & 80 & 92.14 & 92.82 & 网格过粗，攻击下降 \\
$k=4$ & 400 & 92.84 & 99.15 & 最优网格附近 \\
$k=8$ & 80 & 91.94 & 97.55 & 攻击饱和 \\
Seed=0 & 400 & 92.81 & 99.59 & 结果稳定 \\
Seed=1 & 80 & 92.10 & 98.70 & 结果稳定 \\
\bottomrule
\end{longtable}

\section{讨论}

\subsection{为什么 WaNet 难以被传统防御发现}

WaNet 与传统 patch-based 后门不同。传统触发器通常是局部、固定、高对比度图案，防御方法可以通过输入扰动、触发器逆向优化或神经元激活异常发现它们。WaNet 的触发器是全局空间形变，没有明显局部图案，也不一定对应某个简单可视化 mask。因此，依赖局部异常或模型输出熵的防御方法会遇到困难。

Noise Mode 进一步增强了这种隐蔽性。它通过加入随机形变样本，使模型不能简单地把"图像被形变"作为后门信号，而必须学习特定 warping 模式。这会削弱 Neural Cleanse 通过小触发器逆向恢复后门的能力。

\subsection{频域检测的意义与局限}

频域检测的意义在于提供了一个不同视角：不再只问"模型看到触发器后是否表现异常"，而是问"触发器本身是否改变了图像统计"。从本实验看，warping 虽然视觉不可见，但由插值造成的频域变化仍然有检测价值。

但该方法也有明显局限：
\begin{itemize}
    \item 当前 AUC=0.735，仍不足以作为可靠部署级检测器；
    \item 实验主要在 CIFAR-10 和固定 warping 参数下完成，泛化性仍需验证；
    \item 如果攻击者在训练或生成触发器时显式约束频域特征，当前特征可能被绕过；
    \item SVM 使用手工特征，表达能力有限，后续可以考虑学习式频域检测器或多尺度小波特征。
\end{itemize}

因此，频域检测器更适合作为本项目的扩展发现，而不是夸大为已经解决 WaNet 防御问题。更严谨的说法是：我们发现了 WaNet 在图像信号层面的可检测痕迹，并初步验证了该方向优于 STRIP 的模型行为检测。

\section{成员分工}

\begin{itemize}
    \item \textbf{周争妍：}负责论文主复现链路。具体包括：整理 WaNet 论文和官方实现；配置 Conda/GPU 环境；适配 Python、PyTorch、CUDA 版本；编写和维护 tmux 后台运行脚本；完成 MNIST、CIFAR-10、GTSRB 三个数据集在 all2one/all2all 两种模式下的训练与评估；完成 STRIP、Fine-Pruning、Neural Cleanse 三类防御验证；整理 \texttt{eval\_results.json}、防御输出；撰写主复现实验说明。
    \item \textbf{蒋竺君：}负责机制分析和扩展实验。具体包括：(1)~设计并完成 Noise Mode（有/无）、形变强度 $s$（0.25/0.50/1.00）、控制网格大小 $k$（2/4/8）和随机种子（0/1）共 10 组消融训练任务，系统分析各参数对攻击成功率、干净准确率和 Cross Accuracy 的影响；(2)~基于 warping 操作依赖插值运算这一观察，提出并实现基于 FFT 频域特征的 warping 后门检测方案——包括频域特征提取（高频能量比、频谱熵等 6 维特征）、RBF-SVM 分类器训练、ROC/AUC 评估与 STRIP 对比分析；(3)~完成频域可视化工作，包括 FFT 频谱对比图、六维特征分布图、ROC 曲线图的生成与分析；(4)~完成 warping field 变形场可视化、原图与 warped 图对比、PSNR/SSIM 感知质量分析、不同 $(s,k)$ 参数矩阵效果图及消融实验汇总图等可视化；(5)~设计并实现交互式 Web 演示平台（Dashboard），含实时 warping 变形演示、频域检测结果展示和消融实验数据展示，用于答辩现场演示；(6)~制作答辩 PPT 并撰写答辩稿；(7)~撰写本文的消融实验分析、频域检测器及可视化相关章节。
\end{itemize}

\section{伦理说明}

本项目仅在公开数据集和课程实验环境中复现后门攻击，用于理解深度学习模型的安全风险与防御局限。实验目的不是攻击真实系统，而是分析现有防御方法的不足，并探索更有效的检测思路。所有代码、模型和结果仅用于教学与研究。

\section{结论}

本项目复现并分析了 WaNet 后门攻击。实验表明，WaNet 可以在 CIFAR-10 上达到 99.15\% 的攻击成功率，同时保持 92.84\% 的干净准确率，说明基于图像形变的触发器兼具攻击有效性和视觉隐蔽性。防御实验显示，STRIP、Fine-Pruning 与 Neural Cleanse 均难以稳定检测或移除该后门。

消融实验进一步说明，Noise Mode 是 WaNet 绕过防御的关键机制；形变强度 $s$ 和控制网格大小 $k$ 对攻击效果高度敏感，存在较窄的有效参数区间。多种子实验验证了结果的稳定性。最后，提出频域检测器，从图像信号层面检测 warping 痕迹，在 CIFAR-10 上达到 AUC=0.735，高于 STRIP 的随机水平。这说明 WaNet 虽然视觉上隐蔽，但仍可能通过频域统计被部分识别。后续工作可以继续研究更鲁棒的频域/多尺度检测方法，并评估其在更多数据集、模型结构和自适应攻击下的泛化能力。

\begin{thebibliography}{9}

\bibitem{wanet}
Tuan Anh Nguyen and Anh Tuan Tran.
\newblock \textit{WaNet: Imperceptible Warping-based Backdoor Attack}.
\newblock International Conference on Learning Representations, 2021.
\newblock \url{https://openreview.net/forum?id=eEn8KTtJOx}

\bibitem{official}
VinAIResearch.
\newblock \textit{Warping-based Backdoor Attack Official Implementation}.
\newblock \url{https://github.com/VinAIResearch/Warping-based_Backdoor_Attack-release}

\bibitem{project}
axgo-ja.
\newblock \textit{NIS3353 - AI安全课程设计}.
\newblock \url{https://github.com/axgo-ja/NIS3353}

\end{thebibliography}

\end{document}
